%% ICAIF '26 submission — ACM sigconf, double-blind (anonymous).
%% Compile on Overleaf (acmart is built in) or with: pdflatex -> bibtex -> pdflatex x2.
%%
%% ICAIF '26 SUBMISSION VERSION — 8-page HARD limit INCLUDING references, no appendix.
%% This is the conference cut (no appendix). The extended arXiv version with the full
%% appendix is paper.tex. Keep this <= 8 pages when compiled.
%% Remove the `anonymous` option for the camera-ready to reveal the author block below.
\documentclass[sigconf,anonymous,review]{acmart}

\settopmatter{printacmref=false}
\renewcommand\footnotetextcopyrightpermission[1]{}
\pagestyle{plain}
\setcopyright{none}
\acmConference[ICAIF '26]{ACM International Conference on AI in Finance}{November 14--17, 2026}{Milan, Italy}

\usepackage{booktabs}
\usepackage{amsmath}

\begin{document}

\title{Skill or Structure? Auditing Learned Skill and\\ Explanation Faithfulness in Residual Reinforcement Learning\\ for Portfolio Allocation}

%% Author block suppressed by the `anonymous` class option for double-blind review.
\author{Meghanadh}
\affiliation{%
  \institution{Independent Researcher}
  \country{}
}
\email{meghanadh27@gmail.com}
\renewcommand{\shortauthors}{Anonymous}

\begin{abstract}
Reinforcement learning (RL) allocators are often credited with reducing
out-of-sample left-tail risk, yet long-only allocation already inherits downside
protection for free from the simplex constraint, log-reward geometry, and the
classical rebalancing premium. We ask how much of an RL allocator's apparent skill
is genuinely learned rather than inherited structure, and whether post-hoc
explanations of the trained policy identify its true decision mechanism. We train
PPO as a \emph{residual} on a structural-null base policy (equal-weight,
volatility-scaled, or risk-parity) with a \emph{structure-baselined} reward that
credits only return earned above the base, and calibrate the resulting skill
measure on a synthetic market with a toggleable predictive signal: skill rises
monotonically with forecastability and vanishes when prediction is impossible. On
real ETFs (2005--2024, 10\,bps costs, walk-forward), both a scalar de-risking gate
and a more expressive per-asset tilt add \emph{no} skill above structure once
measured net of a phase-randomization placebo null, across three bases and many
seeds; the expressive agent loses more. We then probe the trained policy with
per-feature causal ablation and compare it against gradient saliency and
KernelSHAP. Attribution is broadly faithful to the causal mechanism, but which
method is most faithful is agent-dependent, and SHAP recurrently over-credits a
salient, high-variance feature. The synthetic method-agreement structure transfers
to the real agent, whose decision rule is coherent and faithfully readable but not
skillful.
\end{abstract}

\begin{CCSXML}
<ccs2012>
 <concept><concept_id>10010147.10010257.10010258.10010261.10010269</concept_id>
 <concept_desc>Computing methodologies~Reinforcement learning</concept_desc>
 <concept_significance>500</concept_significance></concept>
 <concept><concept_id>10003752.10010070.10010071</concept_id>
 <concept_desc>Theory of computation~Sequential decision making</concept_desc>
 <concept_significance>300</concept_significance></concept>
 <concept><concept_id>10010147.10010178.10010179</concept_id>
 <concept_desc>Computing methodologies~Machine learning</concept_desc>
 <concept_significance>300</concept_significance></concept>
</ccs2012>
\end{CCSXML}
\ccsdesc[500]{Computing methodologies~Reinforcement learning}
\ccsdesc[300]{Theory of computation~Sequential decision making}

\keywords{reinforcement learning, portfolio allocation, explainability,
attribution, causal intervention, tail risk, residual policy learning}

\maketitle

\section{Introduction}
Model-free reinforcement learning has been reported to beat mean--variance
optimization and reduce out-of-sample left-tail risk in portfolio allocation
\cite{lavko2023,zhang2020dlportfolio}. Two long-standing critiques temper this
claim. First, evaluations frequently omit transaction costs and rarely clear a
strong $1/N$ benchmark \cite{demiguel2009,kruthof2025}. Second, the ``skill''
being measured may not be skill at all: a long-only allocator inherits downside
protection mechanically from the probability-simplex constraint, from maximizing
log wealth, and from the rebalancing premium of Stochastic Portfolio Theory
\cite{fernholz2002}. A volatility-managed or risk-parity rule delivers much of the
reported tail benefit with no forecasting \cite{moreira2017,maillard2010}.

We treat these critiques as a measurement problem and ask two questions.
\textbf{(RQ-skill)} How much of an RL allocator's apparent tail-risk benefit is
learned skill above the free structure, rather than the structure itself?
\textbf{(RQ-faithful)} When we explain the trained policy with standard
attribution tools, do those explanations identify the mechanism the policy
actually uses?

Our approach isolates skill by construction. We fix a \emph{structural-null base
policy} that already harvests the free structure, train the agent as a
\emph{residual} on that base, and reward it only for return earned \emph{above} the
base on the same market path. A zero residual reproduces the base exactly, so the
zero-skill floor is free and the gradient never sees the structural growth. We
calibrate this measure on a synthetic market whose predictive signal can be
switched off, giving a ground-truth boundary where genuine skill is impossible. We
apply the calibrated measure to real ETFs, and finally probe the trained policy's
mechanism with causal interventions, adjudicating gradient saliency and KernelSHAP
against what those interventions reveal.

\noindent\textbf{Contributions.}
(1) A structure-baselined residual-RL skill measure, validated against synthetic
ground truth where skill rises with forecastability and vanishes without it, for a
scalar and an expressive agent.
(2) A real-data verdict: neither a scalar de-risking gate nor an expressive
per-asset tilt adds skill above structure net of a placebo null, across three bases
and many seeds; the expressive agent loses \emph{more}.
(3) A causal-vs-attribution faithfulness probe across three agents, showing
attribution is broadly faithful but its reliability is agent-dependent, with a
recurring SHAP salience bias, and that the synthetic method-agreement transfers to
the real agent. We frame success as a rigorous, honest verdict, not as beating the
market.

\section{Related Work}
RL allocators span value- and policy-based methods and are increasingly paired with
interpretability claims \cite{lavko2023,cong2021alphaportfolio,delarica2025}. Recent
cost- and benchmark-aware replications question whether the RL advantage survives
realistic frictions and a $1/N$ baseline \cite{kruthof2025,demiguel2009}. Our
residual formulation borrows residual policy learning from robotics, where a learned
policy corrects a fixed controller \cite{silver2018,johannink2019}; we repurpose it
as a measurement instrument, pairing it with a structure-baselined reward and a
synthetic ground-truth test. On explanation, gradient saliency \cite{simonyan2014}
and KernelSHAP \cite{lundberg2017} are the standard tools, but saliency maps for
deep RL have been shown to be exploratory rather than explanatory \cite{atrey2020}.
We contribute a controlled, causally grounded test of what these tools recover for
RL allocators specifically.

\section{Method}

\subsection{Structure-baselined residual RL}
Let $b_t \in \Delta^{n-1}$ be the weights of a fixed base policy on the probability
simplex, computed from a trailing window: equal-weight, inverse-volatility
(volatility-scaled), or equal-risk-contribution risk parity \cite{maillard2010}. The
agent acts as a residual on $b_t$ in one of two modes.

\emph{De-risking gate.} A scalar action $g_t \in [0,1]$ shifts weight toward a safe
asset $s$: $w_t = (1-g_t)\,b_t + g_t\,e_s$. Here $g_t=0$ reproduces the base and
larger $g_t$ times de-risking.

\emph{Expressive tilt.} A per-asset action $a_t \in \mathbb{R}^n$ gives
$w_t = \Pi_{\Delta}\big(b_t + \tau\tanh(a_t)\big)$, where $\Pi_\Delta$ is Euclidean
projection onto the simplex \cite{duchi2008} and $\tau$ bounds the tilt. A zero
action again reproduces the base, but the agent can express any long-only deviation.
An early \emph{unbounded} per-asset residual (no $\tanh$/$\tau$ bound) over-tilted
and produced negative skill even synthetically; bounding the tilt fixed this and
motivated reporting both the parsimonious gate and the bounded expressive tilt.

The reward is \emph{structure-baselined}: the agent's net log return minus the
base's net log return on the same path, each charged its own turnover cost,
\begin{equation}
r_t = \log\!\big(1 + w_t^\top \rho_t - c\|w_t - w_{t-1}\|_1\big)
    - \log\!\big(1 + b_t^\top \rho_t - c\|b_t - b_{t-1}\|_1\big),
\end{equation}
where $\rho_t$ are asset returns and $c$ the per-unit cost. The cumulative reward is
realized \emph{skill above structure}; it is zero for the base itself, so the
gradient sees only what the agent adds. We train PPO \cite{schulman2017} with an MLP
policy via Stable-Baselines3 \cite{raffin2021}. The observation is a trailing return
window plus realized volatility, and (for the tilt) dual-horizon volatility,
momentum, and the base weights, with a de-risking signal.

\subsection{Synthetic ground truth}
To calibrate the measure we generate regime-switching markets with a \emph{toggleable}
leading signal of strength $\sigma\in[0,1]$: a two-asset risky/safe world for the
gate and a multi-asset world for the tilt. When $\sigma$ is high, timed de-risking is
possible; when $\sigma=0$ the market is unforecastable and genuine skill is
impossible by construction. A valid measure must rise with $\sigma$ and vanish at
$\sigma=0$.

\subsection{Real-data placebo null}
On real data there is no ground truth for ``no skill.'' We report skill \emph{net of
a phase-randomization placebo null}: surrogate series that preserve each asset's
power spectrum but destroy volatility clustering, retrained identically. The placebo
mean is the luck/overfitting floor; a positive floor is expected and is exactly why
netting is necessary. We report skill net of this null with a small-sample $t$
confidence interval and the fraction of placebo draws that match or beat the real
agent (placebo exceedance). Agents are evaluated with an anchored expanding-window
walk-forward and per-fold retraining, and placed beside $1/N$, minimum-variance, and
CVaR-minimizing baselines on the same out-of-sample window.

\subsection{Causal-vs-attribution faithfulness probe}
Given a trained policy we replay its deterministic decisions and measure two tracks
over the observation features. The \emph{causal} track ablates each feature (freeze
to its mean; permute across time) and records the mean absolute change in the
policy's scalar decision, a direct measure of what the policy responds to; we also
inject volatility shocks and flip the signal as legibility checks. The
\emph{attribution} track computes gradient saliency \cite{simonyan2014} and
KernelSHAP \cite{lundberg2017}.

A scale confound must be removed first. The signal feature has $35$--$100\times$ the
raw scale of the return and volatility features, and raw gradient saliency is
scale-invariant while ablation and SHAP are scale-sensitive; left unfixed, the
comparison would be dominated by units and could \emph{falsely} read as unfaithful.
We therefore express saliency in gradient$\times$standard-deviation units and put all
four methods on one basis of normalized per-feature shares, computing the causal
group verdict from per-feature (not joint) ablation. We report per-feature Spearman
correlation between the causal profile and each attribution profile. Because
normalized shares always sum to one, they hide an inactive agent; we additionally
report \emph{activity diagnostics} (the standard deviation of the policy's decision
over the rollout and the raw, un-normalized causal magnitude per group) and a
degenerate-agent guard. On synthetic data the leading signal is the known driver; on
real data no oracle exists, so only \emph{method agreement} is adjudicable.

\section{Experiments}
The real universe is 11 liquid ETFs spanning equity sectors, Treasuries, gold, and
international equity (2005--2024, daily, 10\,bps costs). Synthetic runs use 150k
training steps and five seeds; real robustness sweeps use five agent seeds and ten
placebo draws per base. All reported figures are held out.

\subsection{The skill measure is valid (RQ-skill, synthetic)}
Table~\ref{tab:rq2} shows mean skill rising monotonically with signal strength for
the gate and vanishing toward the unforecastable null. The expressive tilt cannot
reach a costless do-nothing floor---with no signal it over-churns and loses about
$-6.4\times10^{-5}$---so it is judged net of that churn null: skill-on minus
skill-off is $+2.0\times10^{-4}$, about four times the floor. The measure detects
learned skill only when learned skill is possible, for both agents. A caveat we
carry forward: at higher training budgets the signal-off floor is noisy and slightly
positive, because PPO extracts small \emph{spurious} skill from pure noise; that
floor is the overfitting baseline and is why the real-data analysis nets against a
matched null (\S\ref{sec:rq1}).

\begin{table}[t]
  \caption{Synthetic validation: mean structure-baselined skill (held-out, 5 seeds).
  Skill rises with forecastability and vanishes at the null.}
  \label{tab:rq2}
  \small
  \begin{tabular}{lc}
    \toprule
    Signal strength $\sigma$ (gate) & Mean skill \\
    \midrule
    0.00 (unforecastable null)   & $5.9\times10^{-5}$ \\
    0.50                         & $1.7\times10^{-4}$ \\
    0.95 (strongly forecastable) & $2.9\times10^{-4}$ \\
    \midrule
    \multicolumn{2}{l}{\emph{Expressive tilt}, net-of-null (signal 0.95):}\\
    \quad skill off / skill on            & $-6.4\times10^{-5}$ / $+1.4\times10^{-4}$\\
    \quad skill net of churn floor        & $+2.0\times10^{-4}$ ($4.1\times$)\\
    \bottomrule
  \end{tabular}
\end{table}

\subsection{No skill above structure (RQ-skill, real ETFs)}\label{sec:rq1}
Tables~\ref{tab:rq1gate} and~\ref{tab:rq1tilt} report skill net of the placebo null
across the three bases, for the gate and the tilt. Every skill-net confidence
interval lies entirely below zero, the placebo (luck) mean is positive everywhere
(so netting is essential), and every structureless surrogate beats the real agent
(placebo exceedance $=1.00$) for all bases and both agents. The gate learns to mostly
do nothing (mean gate $\approx0.04$), tracking its base on all tail and risk metrics
while both clearly beat naive $1/N$; the tail benefit over $1/N$ is therefore
inherited structure, not learned skill. Skill-net is most negative for the weakest base
(equal-weight), where the gate activates most and churns without edge. The expressive
tilt loses \emph{more} than the gate ($1.5$--$1.8\times$) in the same base ordering:
more expressiveness means more churn, and where no timeable structure exists that
churn is pure cost drag. A capable agent sharpens the verdict rather than rescuing it.

\begin{table}[t]
  \caption{Real ETFs, \emph{gate}: skill net of the placebo null (5 agent seeds, 10
  placebo draws per base). All intervals below zero; exceedance $=1.00$.}
  \label{tab:rq1gate}
  \small
  \begin{tabular}{lcccc}
    \toprule
    Base & Agent & Placebo & Skill net & 95\% CI ($\times10^{-4}$)\\
    \midrule
    Equal-weight & $-8.6$e-5 & $+9.0$e-5 & $-1.76$e-4 & $[-2.9,-0.7]$\\
    Vol-scaled   & $-7.1$e-5 & $+4.8$e-5 & $-1.19$e-4 & $[-2.1,-0.3]$\\
    Risk-parity  & $-6.0$e-5 & $+4.1$e-5 & $-1.01$e-4 & $[-1.4,-0.6]$\\
    \bottomrule
  \end{tabular}
\end{table}

\begin{table}[t]
  \caption{Real ETFs, \emph{expressive tilt}: skill net of the null. Same sweep;
  all intervals below zero, exceedance $=1.00$. Loses more than the gate.}
  \label{tab:rq1tilt}
  \small
  \begin{tabular}{lccc}
    \toprule
    Base & Agent & Skill net & vs.\ gate\\
    \midrule
    Equal-weight & $-1.56$e-4 & $-2.06$e-4 & $1.17\times$\\
    Vol-scaled   & $-1.26$e-4 & $-1.75$e-4 & $1.47\times$\\
    Risk-parity  & $-1.03$e-4 & $-1.54$e-4 & $1.52\times$\\
    \bottomrule
  \end{tabular}
\end{table}

\subsection{Is the explanation faithful? (RQ-faithful)}
Table~\ref{tab:rq3} summarizes the probe across three agents. On the synthetic gate,
where the signal is the known driver, KernelSHAP names it the top feature in all five
seeds and gradient saliency in four of five---the miss is the least
signal-concentrated agent (causal signal share $0.54$ vs.\ $\approx1.0$)---and
per-feature agreement with the causal profile is positive for both.

On the synthetic expressive tilt, saliency is near-perfectly aligned with the causal
profile ($+0.999$) while SHAP over-credits the salient, high-variance signal feature
($+0.88$), naming it the top group in four of five seeds when the causal analysis
does not; this inverts which method is more reliable relative to the gate. A
methodological caveat appears here: with wildly unequal group cardinalities (100
return features vs.\ 1 signal), group-level ``top driver'' verdicts are
cardinality-confounded---per feature the signal is the strongest single feature
($\sim54\times$ an average return feature)---so the per-feature Spearman is the robust
measure.

On the real agent there is no ground truth, so we report method agreement only. The
activity diagnostic overturns a natural prior: although the mean gate is
$\approx0.05$, its standard deviation is $\approx0.18$, so the agent actively swings
the gate rather than sitting still. It has a real, if unprofitable, decision
mechanism (none of five seeds degenerate). Given that live mechanism, per-feature
agreement with causal ablation is strong---saliency $+0.996$, SHAP $+0.79$---
preserving the saliency-over-SHAP ordering seen synthetically. All four methods
attribute about $0.96$ of the normalized share to the recent-returns block, with the
crisis signal a minor contributor; unlike the synthetic tilt this is not a
cardinality artifact, since the signal is only about three times an average return
feature. The real allocator has a coherent, faithfully readable decision rule; the
rule simply is not skillful.

\begin{table}[t]
  \caption{Faithfulness probe: per-feature Spearman between the causal profile and
  each attribution method (5 seeds). Real-data numbers are method agreement, not
  faithfulness (no ground-truth driver).}
  \label{tab:rq3}
  \small
  \begin{tabular}{lccl}
    \toprule
    Agent & Saliency & SHAP & Note\\
    \midrule
    Synthetic gate & $+0.79$ & $+0.61$ & SHAP names driver 5/5; saliency 4/5\\
    Synthetic tilt & $+0.999$ & $+0.88$ & SHAP over-credits salient signal\\
    Real gate      & $+0.996$ & $+0.79$ & agreement only; agent is active\\
    \bottomrule
  \end{tabular}
\end{table}

\section{Discussion and Limitations}
The questions close a loop. The measure detects skill when it exists (synthetic),
finds none on real markets net of luck (both agents, three bases), and shows the
trained policy's mechanism is largely readable by post-hoc attribution, with saliency
tracking the causal per-feature profile more tightly than SHAP across all three
agents. The recurring practical failure mode is SHAP over-crediting a salient,
high-variance feature; which tool to trust is agent-dependent.

Several limitations bound these claims. The synthetic market is deliberately simple,
chosen so that ``no skill'' is well defined rather than realistic. The real study uses
one liquid universe and daily data; other universes, frequencies, or richer action
spaces may differ. Attribution on real data measures agreement, not faithfulness,
because no oracle mechanism exists; high agreement cannot rule out two methods being
wrong together, though causal ablation is the most direct available reference. The
per-feature causal vectors were aggregated to groups rather than persisted in full; a
camera-ready extension should save them to state per-feature dominance directly.
Finally, the placebo null destroys volatility clustering but preserves linear
autocorrelation; alternative nulls are worth reporting.

\section{Conclusion}
Treating RL portfolio allocation as a measurement problem, we separate learned skill
from inherited structure with a structure-baselined residual and a synthetic
ground-truth calibration, and test whether standard explanations identify the trained
policy's mechanism. On real ETFs, neither a parsimonious nor an expressive RL
allocator earns skill above a transparent structural base once costs and a luck null
are accounted for. The policy's decision rule is nonetheless coherent and faithfully
readable, most reliably by gradient saliency; KernelSHAP carries a salience bias.
Code and configurations will be released publicly.

%% --- References ---
\bibliographystyle{ACM-Reference-Format}
\bibliography{references}

\end{document}
